\documentclass[12pt]{article}
\usepackage[utf8]{inputenc}
\usepackage[T1]{fontenc}
\usepackage{hyperref}
\title{Owa Language Resources for K-12 Curriculum Development}
\author{Compiled by Claude Sonnet 5.5 (claude-sonnet-5-5) for Andr\'as Kornai}
\date{}
\begin{document}
\maketitle

\section{Introduction}
Owa (ISO 639-3 \texttt{stn}, formerly called Santa Ana) is an Austronesian language of the Solomon Islands. It is spoken in the south of Makira and on nearby islands, including Owariki (Santa Catalina). About 8,400 speakers are reported. It is closely related to Kahua.

\section{Primary Educational Resources}
No primers, textbooks or school programmes for Owa turned up in the searches.

\section{Linguistic and Reference Materials}
\begin{itemize}
\item Greg Mellow, \emph{A Dictionary of Owa: A Language of the Solomon Islands} (De Gruyter, Pacific Linguistics, 2013). It has more than 3,900 entries with examples and includes an overview of the phonology, morphology and syntax. Publisher page: \url{https://www.degruyterbrill.com/document/doi/10.1515/9781614513612/html}
\item Library record for the dictionary: \url{https://oeaw.ac.at/resources/Record/9781614513612}
\item Wikipedia, ``Owa language'': \url{https://en.wikipedia.org/wiki/Owa_language}
\item Omniglot, Owa writing and orthography notes: \url{https://omniglot.com/writing/owa.htm}
\end{itemize}

\section{Religious and Community Resources}
Anglican missionaries began writing the language in the 1930s and several religious texts were translated into Owa. The searches did not return specific titles or URLs for them. The Wikipedia article on Owariki describes the island community: \url{https://en.wikipedia.org/wiki/Owariki}

\section{Digital and Audio Resources}
No audio, video or software was found. The searches did not turn up a free online copy of the dictionary, although some sections may be available as PDFs from the publisher.

\section{Curriculum Development Context}
Owa is a stable first language of its community and is not known to be taught in schools. The Mellow dictionary is the main reference for orthography and vocabulary. A curriculum would need new primers and readers, and would need to engage local church and community bodies, whom the searches did not identify.

\end{document}